\chapter{Практика}
\label{ch:intro}

\section{Работа с двум SDR}

\section{Идея}

Добавим возможность при запуске бинарника передавать через консоль параметры для SDR: uri, mode, time, где uri - адрес SDR 
(таким образом мы сможем удобно выбрать конкретную SDR из двух, подключенных к хосту), mode - режим работы SDR (0 - RX, 1 - TX), 
time - время, которое SDR будет работать (секунды!). В программе будем анализировать аргументы и запускать SDR нужным образом. \\

Дополнительно для удобства добавим структуру, которая будет хранить настройки SDR. Прием и передачу разобьем на отдельные функции

\section{Реализация}

Из функций реализуем отдельную библиотеку. Заголовочный файл:

\begin{lstlisting}
#include <SoapySDR/Device.h>
#include <SoapySDR/Formats.h>
#include <complex.h>
#include <stdint.h>
#include <stdio.h>
#include <stdlib.h>

/**
 * @brief Config for init Pluto SDR
 */
struct sdr_config_t
{
  char *usb_uri;   /**< uri address of pluto (use iio_info -a to see connected plutos)*/
  int buffer_size; /**< count of I/Q samples, that SDR receive/send for one time */
  int tx_sample_rate;
  int tx_carrier_freq;
  float tx_gain;
  int rx_sample_rate;
  int rx_carrier_freq;
  float rx_gain;
  int channels[1] = {0};
};

/**
 * @brief create sdr object
 *
 * @param config pluto config
 * @return pluto object
 */
struct SoapySDRDevice *setup_pluto_sdr(sdr_config_t *config);

/**
 * @brief setup streams
 *
 * @param sdr sdr object
 * @param config sdr config
 * @param isRx 0-setup TX stream, 1-setup RX stream
 * @return stream object
 */
struct SoapySDRStream *setup_stream(struct SoapySDRDevice *sdr,
                                    sdr_config_t *config, bool isRx);

/**
 * @brief close streams, free memory. Finished work with pluto
 *
 * @param sdr sdr object
 * @param rxStream
 * @param txStream
 *
 * @return stream object
 */
void close_pluto_sdr(struct SoapySDRDevice *sdr, SoapySDRStream *rxStream,
                     SoapySDRStream *txStream);

/**
 * @brief start receiving samples and write their in file
 *
 * @param sdr sdr object
 * @param rx_stream rx stream object
 * @param rx_buffer how many samples sdr receive for one time
 * @param buff_size how many samples sdr receive for one time
 * @param rx_mtu Max Transmission Unit
 * @param filename file, where will be writing receive samples
 * @param work_time time, that sdr will receive signal
 */
void start_rx(struct SoapySDRDevice *sdr, SoapySDRStream *rxStream,
              int16_t *rx_buffer, int buff_size, int rx_mtu, char *filename,
              int work_time);

/**
 * @brief start transmit samples
 *
 * @param sdr_object sdr object
 * @param tx_stream
 * @param rx_stream
 * @param tx_samples samples that will transmited
 * @param tx_samples_count size of tx samples
 * @param rx_buffer
 * @param tx_buffer
 * @param filename file, where will be writing tx samples
 * @param work_time time, that sdr will tx signal
 */
void start_tx(struct SoapySDRDevice *sdr, SoapySDRStream *txStream,
              SoapySDRStream *rxStream, int16_t *tx_samples,
              int tx_samples_count, int16_t *rx_buffer, int16_t *tx_buffer,
              int buff_size, int work_time);
\end{lstlisting}

Реализацию функций можно посмотреть в src\/main\/pluto\_lib \\

Пример работы с SDR:

\begin{lstlisting}
/*
  This code work with Pluto SDR. He init and setup SDR, receive/transmit samples
*/

#include <SoapySDR/Device.h>
#include <SoapySDR/Formats.h>
#include <complex.h>
#include <iostream>
#include <stdint.h>
#include <stdlib.h>
#include <unistd.h>
#include <vector>

#include "../includes/general/subfuncs.hpp"
#include "../includes/pluto_lib/pluto_lib.hpp"

int main(int argc, char *argv[]) {

  /*check arguments*/
  if (argc != 4) {
    printf("Programm wait 3 argument:\n \
      1.char* pluto usb_uri (example: )\n \
      2.int rx or tx (0-rx, 1-tx)\n \
      3.int time to work (in seconds)");

    return 1;
  }

  /*tx or rx mode*/
  int mode = std::stoi(argv[2]);

  int time_to_work = std::stoi(argv[3]);

  /*create config for SDR*/
  sdr_config_t config;
  config.usb_uri = argv[1];
  config.buff_size = 1920;
  config.rx_carrier_freq = 800e6;
  config.tx_carrier_freq = 800e6;
  config.rx_sample_rate = 1e6;
  config.tx_sample_rate = 1e6;
  config.rx_gain = 35.0;
  config.tx_gain = -15.0;

  /*files for tx/rx samples*/
  char *rxdata = "../pcm/rxdata.pcm";
  char *txdata = "../pcm/txdata.pcm";

  /*setup pluto*/
  struct SoapySDRDevice *sdr = setup_pluto_sdr(&config);
  struct SoapySDRStream *rxStream = setup_stream(sdr, &config, 1);
  struct SoapySDRStream *txStream = setup_stream(sdr, &config, 0);

  /*create buffers*/
  int16_t tx_buffer[2 * config.buff_size];
  int16_t rx_buffer[2 * config.buff_size];

  /*if we want only rx*/
  if (mode == 0) {
    start_rx(sdr, rxStream, rx_buffer, config.buff_size, config.buff_size,
             rxdata, time_to_work);
  /*if we want only tx*/
  } else if (mode == 1) {
    /*read sample as std::complex*/
    std::vector<std::complex<int16_t>> complex_samples =
        read_pcm(std::string("../pcm/tx_samples.pcm"));
    /*create array in format: I0,Q0,I1,Q1 ... In,Qn*/
    int uncomplex_samples_size = complex_samples.size() * 2;
    std::vector<int16_t> uncomplex_samples;
    uncomplex_samples.reserve(uncomplex_samples_size);

    for (int i = 0; i < complex_samples.size(); ++i) {
      uncomplex_samples.push_back(complex_samples[i].real());
      uncomplex_samples.push_back(complex_samples[i].imag());
    }

    /*start tx*/
    start_tx(sdr, txStream, rxStream, rx_buffer, uncomplex_samples.data(),
             complex_samples.size(), config.buff_size, txdata, time_to_work);
  } else {
    printf("Invalid mode. Enter mode=0 or mode=1");
  }

  /*free resources*/
  close_pluto_sdr(sdr, rxStream, txStream);

  return EXIT_SUCCESS;
}
\end{lstlisting}

\section{Тестирование схемы Гарднера}

Для проверки схемы Гарднера я отправлю некоторый текст, а потом сделаю прием без схемы Гарднера и со схемой Гарднера \\

Отправляемый текст:

\begin{figure}[H]
    \centering
    \includegraphics[width=1.0\textwidth]{text.png}
    \caption{Отправляемый текст}
\end{figure}

Без Гарднера:

\begin{figure}[H]
    \centering
    \includegraphics[width=1.0\textwidth]{wog.png}
    \caption{Принятое сообщение (без Гарднера)}
\end{figure}

Видим полную "кашу" из символов. Ничего общего это сообщение не имеет с отправляемым \\

С Гарднером

\begin{figure}[H]
    \centering
    \includegraphics[width=1.0\textwidth]{wg.png}
    \caption{Приянтое сообщение (с Гарднером)}
\end{figure}

Сообщение не идеальное, но уже очень похоже на отправляемое. Результат очень зависит от параметров самой схемы Гарднера и нужно правильно ее настроить.

\endinput
